problem:  
Let \(S^{n+1}\) be the round sphere (\(n \ge 7\)), and let \(\mathcal{M}_n\) be the space of smooth Riemannian metrics on \(S^{n+1}\) with the \(C^3\)-topology.  
For each \(g \in \mathcal{M}_n\), denote by \(\mathcal{A}(g)\) the set of area‑minimizing boundaries mod 2 in \((S^{n+1}, g)\) in the sense of Almgren–Federer.

Suppose \(\Sigma = \partial E \in \mathcal{A}(g)\) and \(x \in \mathrm{Sing}(\Sigma)\) is a singular point **with a unique tangent cone** \(C_x\Sigma\).  
Let its link \(\Gamma_x = C_x\Sigma \cap S_x^n \subset T_xS^{n+1}\) be regarded as a singular minimal hypersurface in the unit round sphere \(S_x^n\).  
On the regular part \(\Gamma_x^{\mathrm{reg}}\), define the classical Jacobi operator
\[
L_{\Gamma_x} = -\Delta_{\Gamma_x} - (|A_{\Gamma_x}|^2 + n - 1)
\]
and assume \(L_{\Gamma_x}\) admits a **self‑adjoint Friedrichs extension** \(L_{\Gamma_x}^{\mathrm{Fr}}\) with **compact resolvent** (hence discrete real spectrum, ordered as \( \lambda_1 < \lambda_2 \le \lambda_3 \le \cdots \)).  

For such \(g, \Sigma, x,\) fix the lowest nontrivial eigenvalue \(\lambda_1(g, \Sigma, x) := \lambda_1(L_{\Gamma_x}^{\mathrm{Fr}})\).  

> **Problem:**  
> Does there exist a comeager subset \(\mathcal{G} \subset \mathcal{M}_n\) such that for every \(g \in \mathcal{G}\),
> whenever \(\Sigma = \partial E \in \mathcal{A}(g)\) and \(x \in \mathrm{Sing}(\Sigma)\) admits a unique tangent cone with spectral compactness as above, the eigenvalue \(\lambda_1(g, \Sigma, x)\) is **nondegenerate** under perturbation of \(g\)?  
>  
> Equivalently: Are zero‑crossings or spectral bifurcations of the Jacobi operator on links of unique tangent cones a **nongeneric phenomenon** in the space of smooth metrics on \(S^{n+1}\)?  
>  
> More formally, is the set of metrics \(g\) for which some such \(\lambda_1(g, \Sigma, x)\) equals a given value (e.g. 0) a meagre subset of \(\mathcal{M}_n\)?

---

Why is it a "good" problem:  

1. **Mathematical soundness and precision:**  
   The spectral objects are carefully defined—the Jacobi operator’s Friedrichs extension and compact resolvent assumption ensure discrete spectrum and that the smallest eigenvalue \(\lambda_1\) is well-defined. This resolves difficulties in definitional rigour often arising from singular links.

2. **Conceptual depth:**  
   The problem investigates **spectral stability of singular tangent cones** under generic metric perturbations, a precise singular analogue of the classical bumpy metric theorem. It explores whether the geometric degeneracies of singular models correspond to nongeneric “resonances” of elliptic operators on singular minimal links.

3. **Bridge between areas:**  
   The question connects:
   - **Geometric measure theory** (tangent cone analysis, area-minimizing singularities);
   - **Spectral theory and functional analysis** (Friedrichs extension and spectral stability);
   - **Global differential topology** (generic properties of metrics and transversality).  
   Addressing it would likely require insights across these domains, fostering cross-pollination between variational geometry and analysis on singular spaces.

4. **Extreme simplicity of statement:**  
   Conceptually, it asks:  
   *“For generic metrics on the sphere, do singular minimizers avoid Jacobi spectral degeneracy?”*  
   The question is simple to articulate but encapsulates profound analytic and geometric intricacies concerning the nature and stability of minimal singularities.

5. **Potential significance:**  
   - A **positive answer** would extend the philosophy of bumpy metrics into singular geometric measure theory, revealing that spectral degeneracy (and hence structural instability) of singular tangent cones is nongeneric.  
   - A **negative answer** would expose intrinsic analytic rigidity of singular minimizing cones that persists regardless of ambient perturbations.  
   Either outcome would redefine understanding of generic regularity phenomena and the analytic landscape of singular variational problems.

In summary, the problem is precise, logically well-founded, and deeply interweaves spectral geometry, metric genericity, and singular minimal surface theory—qualifying unmistakably as a “good” research-level mathematical problem.